\bookchapter{Annexe D : Où trouver des ressources fiables ?}

\booksection{Le support officiel d'abord :}

\begin{itemize}
  \item \textbf{Utilisateurs Apple}: apple.com/cn et iCloud.com.
  \item \textbf{Utilisateurs Android}:Les sites officiels des marques de smartphones ont tous une section « Service et support ».
  \item \textbf{Outils généraux}: Apple utilise iTunes, Huawei l’assistant téléphone Huawei, Xiaomi l’outil de transfert Xiaomi.
\end{itemize}

\booksection{Sites utiles et blogueurs :}

\begin{itemize}
  \item \textbf{Chercher des astuces}:Cherchez sur Bilibili « astuces [votre modèle de smartphone] ».
  \item \textbf{Voir les tests}:Regardez les vidéos de test de plusieurs blogueurs différents.
  \item \textbf{Apprendre la photo}:Suivez des blogueurs de tutoriels de photographie sur smartphone.
\end{itemize}

\booksection{Liste de vérification de sécurité en une phrase :}

\begin{enumerate}
  \item Avez-vous défini un mot de passe d'écran de verrouillage ?
  \item Avez-vous activé « Trouver mon smartphone » ?
  \item Avez-vous activé l'authentification à deux facteurs sur les comptes importants ?
  \item Faites-vous régulièrement des sauvegardes de vos photos et historiques WeChat ?
  \item Retenez la formule anti-arnaque : ne pas croire aveuglément, ne pas divulguer, ne pas transférer d'argent, ne pas être cupide, ne pas paniquer.
\end{enumerate}
